\subsubsection{版本、测量和报告的关联}
源文件冻结清单确定本次提交的 RTL、固件与检查输入；源码标识关联仿真和综合配置；网表哈希关联面积、时序与功耗；窗口编号关联周期与活动数据。报告图表只从这些记录进行加总、单位转换及格式整理，不以重新绘图改变原测试结论。

\begin{figure}[H]\centering
\begin{tikzpicture}[node distance=5mm]
\node[box,text width=119mm] (a) {源码提交与冻结清单：169 个源及证据文件};
\node[box,text width=119mm,below=of a] (b) {同版 RTL、固件、文件顺序、参数及映射网表};
\node[box,text width=119mm,below=of b] (c) {功能结果、周期、面积、时序、功耗及原始截图};
\node[box,text width=119mm,below=of c] (d) {报告数据核算、表格生成、矢量制图及局部截图};
\foreach\u/\v in {a/b,b/c,c/d}\draw[arr](\u)--(\v);
\end{tikzpicture}
\caption{报告与冻结工程的证据关联。最终数值没有混用后续候选的源码或其他版本的网表。}
\end{figure}

报告附带的早期观测快照记录当时 Formality 尚在运行；后续功耗和最终汇总保存已完成的工具输出。两类记录按时间与状态区分，不能用早期状态覆盖后续结论，也不能删除早期失败或告警来使证据表面一致。正文采用后续完成数值，同时明确未匹配点、未约束端点和活动警告仍然存在。

\subsubsection{结果复核顺序}
\begin{enumerate}
\item 校验冻结清单，确认 RTL、固件、编译顺序与候选参数未变化。文件完整性检查不调用 EDA，不等价于重新完成仿真或综合。
\item 读取本地任务和厂家模型任务记录，分别核对输出元素数、写命令数、完成边界及周期。两组任务分开统计。
\item 以同一映射网表核对面积分项、宏实例数和时序约束，保留未约束端点报告；不能只截取 slack 数值。
\item 将十二个活动窗口与功耗逐项配对，检查时间和功耗单位，再计算任务能量、加权平均功耗及均值。
\item 编译报告并逐页检查图表、目录、引用和截图边界；核对 PDF 与源文件包中的数据一致。
\end{enumerate}

图中的结构连接、缓冲生命周期及协议波形均为根据冻结源码绘制的说明图；逐任务曲线和表格来自记录数据；五张测试报告图片保留原文件，只在排版时截取结果区域。三类图使用不同图注说明来源，避免把协议示意误认为实测波形。原始截图中的警告和未完成项随结果一同保留，裁剪只移除无关桌面区域。
